\begin{center}
    \chapter{Analysis and Design}
\end{center}
\label{chap:analysisdesign}

\section*{Introduction}

The previous chapters introduced the host organisation, motivated the
project and surveyed the state of the art on which HireFlow is built.
The present chapter turns this groundwork into a concrete engineering
plan. It first captures the functional and non-functional requirements
through actor identification, a global use-case diagram and a quality
matrix. It then frames the work in a Scrum cadence with a product
backlog and a six-sprint roadmap. Finally, it lays out the physical,
logical and multi-agent architectures of the platform together with
the hardware and software environment that supports its delivery.

\section{Requirements Analysis}

This section enumerates who the platform serves, what it must do and
under which quality constraints it must operate.

\subsection{Identification of the Different Actors}

Two categories of human actor interact with HireFlow, complemented
by a handful of external technical actors whose behaviour the system
has to integrate with.
\begin{itemize}
    \item \textbf{Recruiter} (primary): opens vacancies, drafts and
    validates job descriptions, triggers candidate matching, launches
    interviews, reviews scoring reports and refines the personalisation
    signal. Most of the platform's value is delivered to this actor.
    \item \textbf{Candidate} (primary): applies to a vacancy, uploads a
    \acs{CV}, takes the conversational interview and consults the
    outcome of the application. Interactions are designed to be
    transparent and explainable.
    \item \textbf{External \acs{AI} services} (supporting): the vLLM
    server, the Groq inference endpoint, and the speech models
    (Kokoro \acs{TTS}, Moonshine \acs{ASR}) participate as
    supporting actors invoked synchronously by the backend during the
    interview and scoring flows.
\end{itemize}

\subsection{Functional Requirements}

Functional requirements are organised by epic; each line is uniquely
identified so that traceability with the use-case model, the product
backlog and the test plan can be enforced.

\begin{table}[!htbp]
    \centering
    {\renewcommand{\arraystretch}{0.95}%
    \resizebox{\textwidth}{!}{%
    \begin{tabular}{|l|l|p{9cm}|}
        \hline
        \textbf{ID} & \textbf{Epic} & \textbf{Requirement} \\
        \hline
        FR-01 & JD authoring & Generate a culture-aware JD from a short brief. \\
        FR-02 & JD authoring & Retrieve a relevant subgraph from the talent knowledge graph (Graph RAG). \\
        FR-03 & JD authoring & Allow the recruiter to edit and validate the draft. \\
        \hline
        FR-04 & Matching & Parse uploaded CVs into a structured skill profile. \\
        FR-05 & Matching & Rank candidates against a JD using a deep \acs{GNN} matcher. \\
        FR-06 & Matching & Display per-candidate evidence supporting each rank. \\
        \hline
        FR-07 & Interview & Conduct a voice-driven adaptive interview grounded in the CV. \\
        FR-08 & Interview & Render a Three.js avatar with lip-synced speech. \\
        FR-09 & Interview & Choose between follow-up and topic-switch based on answer depth. \\
        \hline
        FR-10 & Behavioural analysis & Capture face, pose and audio features in real time. \\
        FR-11 & Behavioural analysis & Fuse multimodal features into the four behavioural signals. \\
        \hline
        FR-12 & Scoring & Produce a per-dimension scoring report with an explanation paragraph. \\
        FR-13 & Scoring & Export the report as a downloadable PDF. \\
        \hline
        FR-14 & Personalisation & Learn a per-recruiter taste model from historical decisions. \\
        FR-15 & Personalisation & Surface the taste signal alongside the central recommendation. \\
        \hline
    \end{tabular}%
    }}
    \caption{Functional requirements traced by epic.}
    \label{tab:functional-requirements}
\end{table}

\subsection{General Use Case Diagram}

The diagram in Figure~\ref{fig:global-usecase} summarises the actors
and the high-level capabilities they exercise, together with the
behavioural-signal capture that is factored out as an included
use case.

\begin{figure}[H]
    \centering
    \makebox[\textwidth][c]{\resizebox{1.32\textwidth}{!}{%
    \begin{tikzpicture}[
        font=\small,
        usecase/.style={ellipse, draw=black!70, thick, fill=blue!8,
            minimum width=3.6cm, minimum height=0.9cm, align=center,
            font=\small, inner sep=3pt},
        actorname/.style={align=center, font=\small\bfseries},
        link/.style={thick, black!70},
        dashedlink/.style={thick, dashed, black!60}
    ]
        % System boundary
        \draw[thick, black!50, rounded corners=4pt]
            (-7.2, -5.8) rectangle (7.2, 5.2);
        \node[font=\bfseries\large, fill=white, inner xsep=12pt,
              inner ysep=2pt]
            at (0, 5.2) {HireFlow};

        % Recruiter (left)
        \begin{scope}[shift={(-8.8, 1.0)}]
            \draw[thick] (0, 1.0) circle (0.3);
            \draw[thick] (0, 0.7) -- (0, -0.3);
            \draw[thick] (-0.55, 0.35) -- (0.55, 0.35);
            \draw[thick] (-0.45, -1.0) -- (0, -0.3) -- (0.45, -1.0);
            \node[actorname] at (0, -1.45) {Recruiter};
        \end{scope}

        % Candidate (right)
        \begin{scope}[shift={(8.8, 1.0)}]
            \draw[thick] (0, 1.0) circle (0.3);
            \draw[thick] (0, 0.7) -- (0, -0.3);
            \draw[thick] (-0.55, 0.35) -- (0.55, 0.35);
            \draw[thick] (-0.45, -1.0) -- (0, -0.3) -- (0.45, -1.0);
            \node[actorname] at (0, -1.45) {Candidate};
        \end{scope}

        % Recruiter use cases
        \node[usecase] (uc1) at (-3.8, 4.0)  {Author Job\\Description};
        \node[usecase] (uc2) at (-3.8, 2.3)  {Validate\\Shortlist};
        \node[usecase] (uc3) at (-3.8, 0.6)  {Launch\\Interview};
        \node[usecase] (uc4) at (-3.8, -1.1) {Review\\Scoring Report};
        \node[usecase] (uc5) at (-3.8, -2.8) {Adjust\\Personalisation};

        % Candidate use cases
        \node[usecase] (uc6) at (3.8, 4.0)  {Submit\\Application};
        \node[usecase] (uc7) at (3.8, 2.3)  {Upload CV};
        \node[usecase] (uc8) at (3.8, 0.6)  {Take Interview};
        \node[usecase] (uc9) at (3.8, -1.1) {Consult Decision};

        % Shared / included use case
        \node[usecase, fill=green!10] (uc10) at (0, -4.4)
            {Stream Behavioural\\Signals};

        % Recruiter links
        \draw[link] (-8.15, 1.0) -- (uc1.west);
        \draw[link] (-8.15, 1.0) -- (uc2.west);
        \draw[link] (-8.15, 1.0) -- (uc3.west);
        \draw[link] (-8.15, 1.0) -- (uc4.west);
        \draw[link] (-8.15, 1.0) -- (uc5.west);

        % Candidate links
        \draw[link] (8.15, 1.0) -- (uc6.east);
        \draw[link] (8.15, 1.0) -- (uc7.east);
        \draw[link] (8.15, 1.0) -- (uc8.east);
        \draw[link] (8.15, 1.0) -- (uc9.east);

        % Include relationships — explicit Bezier paths routed through the
        % empty central gap between the two UC columns, with horizontal
        % labels anchored in that gap so they remain readable.
        \draw[dashedlink, -{Latex}] (uc3.south)
            .. controls (-1.4, -1.2) and (-1.4, -3.0) ..
            (uc10.north west);
        \node[font=\scriptsize\itshape, fill=white, inner sep=2pt,
              draw=black!30, rounded corners=1.5pt]
              at (-1.6, -1.2) {«include»};
        \draw[dashedlink, -{Latex}] (uc8.south)
            .. controls ( 1.4, -1.2) and ( 1.4, -3.0) ..
            (uc10.north east);
        \node[font=\scriptsize\itshape, fill=white, inner sep=2pt,
              draw=black!30, rounded corners=1.5pt]
              at ( 1.6, -2.8) {«include»};
    \end{tikzpicture}%
    }}
    \caption{General use-case diagram of HireFlow.}
    \label{fig:global-usecase}
\end{figure}

The diagram makes the responsibility split visible. The recruiter
drives every authoring, evaluation and personalisation activity,
while the candidate is only exposed to the surfaces strictly needed
to apply, interview and consult the outcome. The
\textit{«include»} relationships factor out the behavioural-signal
capture that both \textit{Launch Interview} and \textit{Take
Interview} trigger.

\subsection{Non-Functional Requirements}

The non-functional requirements derive from the dual nature of the
platform: a research-grade multi-agent system that must nevertheless
behave as a usable product.
\begin{itemize}
    \item \textbf{Performance}: end-to-end conversational latency below
    two seconds per turn; behavioural-signal pipeline running at no
    less than fifteen frames per second on a single \acs{GPU}.
    \item \textbf{Scalability}: stateless agent contracts that allow
    additional model instances to be added behind a load balancer
    without code changes.
    \item \textbf{Reliability}: graceful degradation if a single agent
    fails --- the recruiter must always reach a usable shortlist even
    when the conversational layer is unavailable.
    \item \textbf{Observability}: every agent emits structured logs
    capturing the inputs, the version of the model that produced the
    decision and the timestamp at which it was emitted.
    \item \textbf{Usability}: a recruiter with no \acs{AI} background
    must be able to author a \acs{JD}, launch an interview and read
    the scoring report without supervision.
    \item \textbf{Maintainability}: each agent is a self-contained
    Python package with its own tests and contract so that components
    can evolve independently.
    \item \textbf{Portability}: the platform runs on a developer laptop
    and on a single \acs{GPU} node, without any external orchestrator
    dependency.
\end{itemize}

\section{Project Planning with Scrum}

HireFlow was delivered following a Scrum cadence adapted to the
academic context of an end-of-studies project. The methodology was
chosen because it accommodates the exploratory nature of multi-agent
\acs{AI} work --- the precise shape of each agent emerges as data and
models are explored --- while still imposing the discipline of
two- to three-week sprints, a written backlog and a demonstrable
increment at the end of every iteration.

\subsection{Definition of Roles}

The project assigned the three canonical Scrum roles to the actors
of the internship, with each role mapped to a concrete responsibility
inside the team.

\begin{table}[!htbp]
    \centering
    {\renewcommand{\arraystretch}{1.1}%
    \begin{tabular}{|l|p{10.5cm}|}
        \hline
        \textbf{Role} & \textbf{Responsibility and assignment} \\
        \hline
        Product Owner & Owns the product vision, prioritises the backlog
        and validates each sprint demo. Assigned to the Talan Tunisie
        Consulting professional supervisor. \\
        \hline
        Scrum Master & Facilitates the ceremonies, removes impediments
        and protects the team's focus. Assigned jointly to the academic
        supervisor at \acs{ENSI} and the professional mentor at Talan. \\
        \hline
        Development Team & Designs, implements and validates each
        increment. Composed of the \acs{PFE} student acting as a
        single full-stack \acs{AI} developer. \\
        \hline
    \end{tabular}}
    \caption{Scrum roles and their assignment.}
    \label{tab:scrum-roles}
\end{table}

Three ceremonies were observed throughout the internship: a weekly
sprint review with the Product Owner, a fortnightly sprint planning
session, and a daily written stand-up logged in the project tracker.
The Definition of Done required, for every user story, working code
merged on the main branch, unit tests passing in continuous
integration, a short demo recording and a paragraph appended to the
engineering log.

\subsection{Product Backlog}

The product backlog organises the functional requirements of
Section~3.1.2 into user stories. Each story follows the canonical
\textit{As a~/ I want~/ so that} template and carries a \acs{MoSCoW}
priority. Story points are estimated on a Fibonacci scale and
indicate complexity rather than calendar effort.

\begin{table}[!htbp]
    \centering
    {\renewcommand{\arraystretch}{0.95}%
    \resizebox{\textwidth}{!}{%
    \begin{tabular}{|l|l|p{9.5cm}|c|c|}
        \hline
        \textbf{ID} & \textbf{Epic} & \textbf{User story} & \textbf{Pts} & \textbf{Prio.} \\
        \hline
        US-01 & JD authoring & As a recruiter, I want to generate a culture-aware JD from a short brief so that I save authoring time.                                & 8 & M \\
        US-02 & JD authoring & As a recruiter, I want the agent to consult a talent knowledge graph (Graph RAG) so that the JD reflects the company's skill landscape. & 5 & M \\
        US-03 & JD authoring & As a recruiter, I want to edit and validate the draft so that I keep control on the final wording.                                      & 2 & M \\
        \hline
        US-04 & Matching     & As a recruiter, I want uploaded CVs parsed into a structured skill profile so that matching becomes possible.                            & 5 & M \\
        US-05 & Matching     & As a recruiter, I want candidates ranked against a JD using a deep GNN matcher so that I focus on the most relevant profiles.            & 8 & M \\
        US-06 & Matching     & As a recruiter, I want the evidence behind each rank so that I can trust and defend the recommendation.                                  & 3 & M \\
        \hline
        US-07 & Interview    & As a candidate, I want a voice-driven adaptive interview so that the conversation feels natural.                                         & 8 & M \\
        US-08 & Interview    & As a candidate, I want to see a Three.js avatar that lip-syncs the questions so that the exchange is engaging.                           & 5 & S \\
        US-09 & Interview    & As a recruiter, I want the agent to decide between follow-up and topic-switch so that depth is preserved.                                & 5 & M \\
        \hline
        US-10 & Behavioural  & As a recruiter, I want face, pose and audio features captured in real time so that behavioural signals can be computed.                  & 5 & M \\
        US-11 & Behavioural  & As a recruiter, I want the four behavioural signals fused into a profile so that the scoring report is enriched.                         & 5 & M \\
        \hline
        US-12 & Scoring      & As a recruiter, I want a per-dimension scoring report with explanations so that I understand each verdict.                               & 5 & M \\
        US-13 & Scoring      & As a recruiter, I want to export the report as a PDF so that I can share it with the hiring panel.                                       & 2 & S \\
        \hline
        US-14 & Personalisation & As a recruiter, I want a private taste model learned from my decisions so that suggestions adapt to my style.                         & 8 & C \\
        US-15 & Personalisation & As a recruiter, I want the taste signal displayed next to the central rank so that I keep oversight.                                  & 3 & C \\
        \hline
    \end{tabular}%
    }}
    \caption{Product backlog (M = Must, S = Should, C = Could).}
    \label{tab:backlog}
\end{table}

\subsection{Sprint Planning}

The fifteen stories were grouped into six sprints by thematic
proximity rather than strict priority order, so that each iteration
delivered a self-contained increment that the Product Owner could
demonstrate end-to-end. Two of the sprints --- Sprint~1 and Sprint~4
--- carry an explicit \acs{CRISP-DM} cycle because they involve a
modelling experiment whose result feeds the rest of the platform.

\begin{table}[!htbp]
    \centering
    {\renewcommand{\arraystretch}{1.1}%
    \resizebox{\textwidth}{!}{%
    \begin{tabular}{|l|p{4.0cm}|c|p{4.4cm}|p{2.6cm}|}
        \hline
        \textbf{\#} & \textbf{Theme} & \textbf{Dur.} & \textbf{Stories} & \textbf{Methodology} \\
        \hline
        S1 & LLM fine-tuning for the recruiter persona & 3w & US-07, US-09           & Scrum + CRISP-DM \\
        \hline
        S2 & Conversational interview                        & 3w & US-07, US-08, US-09    & Scrum \\
        \hline
        S3 & Computer-vision component and avatar            & 3w & US-08, US-10, US-11    & Scrum \\
        \hline
        S4 & Job matching                                    & 3w & US-04, US-05, US-06    & Scrum + CRISP-DM \\
        \hline
        S5 & Culture-aware job description                   & 2w & US-01, US-02, US-03    & Scrum \\
        \hline
        S6 & Taste module                                    & 2w & US-14, US-15           & Scrum \\
        \hline
    \end{tabular}%
    }}
    \caption{Six-sprint roadmap of the HireFlow internship.}
    \label{tab:sprints}
\end{table}

The roadmap is deliberately bottom-up. The persona-fine-tuned
\acs{LLM} produced in Sprint~1 underpins both the conversational
interview of Sprint~2 and the culture-aware \acs{JD} authoring of
Sprint~5; the matching back-end built in Sprint~4 is reused by the
personalisation layer of Sprint~6. This ordering keeps the demos
incremental and lets each sprint be evaluated against a working
baseline rather than against a moving target.

\section{Application Architecture}

The platform is built as a single deployable unit that nevertheless
respects three classical architectural views. Section~3.3.1 fixes
the \textbf{physical} topology --- which artefacts run on which
machine, with which protocols. Section~3.3.2 then reads the same
system through the \textbf{logical} lens of a three-tier web
application: presentation, business, data. Section~3.3.3 zooms into
the business tier to expose the \textbf{multi-agent} composition
that gives HireFlow its specificity.

\subsection{Physical Architecture}

The deployment is intentionally compact: one developer client, one
\acs{GPU}-equipped server, and one external inference provider for
fall-back generation. This keeps the platform reproducible on a
single Lightning AI Studio node, which was the constraint imposed
by the academic environment.

\begin{figure}[H]
    \centering
    \makebox[\textwidth][c]{\resizebox{1.12\textwidth}{!}{%
    \begin{tikzpicture}[
        font=\footnotesize,
        tier/.style={rectangle, draw=black!75, very thick, rounded corners=6pt,
            fill=#1, inner sep=10pt},
        title/.style={font=\footnotesize\bfseries, anchor=north},
        comp/.style={rectangle, draw=black!60, thick, rounded corners=2pt,
            fill=white, minimum width=3.6cm, minimum height=0.75cm,
            align=center, font=\scriptsize, inner sep=3pt},
        compw/.style={rectangle, draw=black!60, thick, rounded corners=2pt,
            fill=white, minimum width=8.4cm, minimum height=0.75cm,
            align=center, font=\scriptsize, inner sep=3pt},
        sub/.style={rectangle, draw=black!60, thick, rounded corners=2pt,
            fill=white, minimum width=2.5cm, minimum height=0.95cm,
            align=center, font=\scriptsize, inner sep=2pt},
        link/.style={thick, black!75, -{Latex[length=2.4mm]}},
        dashlink/.style={thick, black!70, -{Latex[length=2.4mm]}, dashed}
    ]
        % --- Client tier container ---
        \draw[tier={blue!5}] (-9.6, 0.0) rectangle (-5.0, 5.6);
        \node[title] at (-7.3, 5.5) {Client tier};
        \node[comp] at (-7.3, 4.55) {Chromium browser};
        \node[comp] at (-7.3, 3.65) {React 18 SPA (Vite)};
        \node[comp] at (-7.3, 2.75) {Three.js avatar (R3F)};
        \node[comp] at (-7.3, 1.85) {WebRTC microphone};
        \node[comp] at (-7.3, 0.95) {Scoring-report viewer};

        % --- Server tier container ---
        \draw[tier={orange!8}] (-2.6, 0.0) rectangle (9.4, 5.6);
        \node[title] at (3.4, 5.5) {Server tier (GPU node)};
        \node[compw] at (3.4, 4.55) {FastAPI gateway (Uvicorn)};
        \node[compw] at (3.4, 3.55) {LangGraph orchestrator + A2A clients};
        \node[sub]   at (-1.0, 2.05) {SQLite\\hireflow.db};
        \node[sub]   at ( 1.9, 2.05) {vLLM};
        \node[sub]   at ( 4.9, 2.05) {Kokoro TTS\\+ Moonshine};
        \node[sub]   at ( 7.9, 2.05) {Redis\\(A2A streams)};
        \node[sub, minimum width=11.6cm] at ( 3.4, 0.85) {Local file directory (CVs / PDF reports)};

        % --- External tier container ---
        \draw[tier={green!7}] (0.4, -3.4) rectangle (6.4, -0.8);
        \node[title] at (3.4, -0.9) {External provider};
        \node[comp] at (3.4, -2.35) {Groq Cloud\\LLM fall-back};

        % Inter-tier links
        \draw[link]     (-5.0, 2.8) -- node[above, font=\scriptsize] {HTTPS / WSS} (-2.6, 2.8);
        \draw[dashlink] (3.4, 0.0)  -- node[right, font=\scriptsize] {fall-back} (3.4, -0.8);
    \end{tikzpicture}%
    }}
    \caption{Physical deployment architecture of HireFlow.}
    \label{fig:physical-arch}
\end{figure}

The client requires nothing more than a recent browser supporting
WebGL and WebRTC. All trained weights, the SQLite store and the
object directory live on the server, which means that no candidate
data ever leaves the perimeter under normal operation. The dashed
arrow towards Groq represents a fall-back path used only when the
local vLLM engine is unavailable; that path is gated behind a
feature flag so it can be disabled when stricter data-residency
constraints apply.

\subsection{System Logical Architecture}

Mapped onto the classical three-tier model, the platform's
responsibilities partition cleanly into a presentation tier, a
business tier and a data tier. The business tier is itself the
locus of the multi-agent system that Section~3.3.3 details.

\begin{figure}[H]
    \centering
    \makebox[\textwidth][c]{\resizebox{1.12\textwidth}{!}{%
    \begin{tikzpicture}[
        font=\footnotesize,
        band/.style={rectangle, draw=black!75, very thick, rounded corners=4pt,
            minimum width=15.2cm, minimum height=#1, inner sep=8pt},
        sideband/.style={rectangle, draw=black!75, very thick, rounded corners=4pt,
            fill=black!10, minimum width=2.4cm, minimum height=#1, align=center,
            font=\footnotesize\bfseries, inner sep=4pt},
        comp/.style={rectangle, draw=black!60, thick, fill=white,
            rounded corners=2pt, minimum width=3.0cm, minimum height=0.95cm,
            align=center, font=\scriptsize, inner sep=3pt},
        link/.style={thick, black!75, -{Latex[length=2.6mm]}}
    ]
        % --- Presentation tier ---
        \node[sideband=1.6cm] (preslbl) at (-9.0, 4.0)
            {Presentation\\tier};
        \node[band=1.6cm, fill=blue!5] (pres) at (0.5, 4.0) {};
        \node[comp] at (-5.5, 4.0) {React 18 SPA\\(Vite, Tailwind)};
        \node[comp] at (-1.8, 4.0) {Three.js + R3F\\avatar};
        \node[comp] at ( 1.8, 4.0) {WebRTC\\audio stream};
        \node[comp] at ( 5.5, 4.0) {Scoring-report\\viewer};

        % --- Business tier ---
        \node[sideband=3.6cm] (bizlbl) at (-9.0, 0.5)
            {Business\\tier};
        \node[band=3.6cm, fill=orange!8] (biz) at (0.5, 0.5) {};
        \node[comp] at (-5.5, 1.4) {FastAPI\\gateway};
        \node[comp] at (-1.8, 1.4) {LangGraph\\orchestrator};
        \node[comp] at ( 1.8, 1.4) {A2A clients\\(JSON-RPC)};
        \node[comp] at ( 5.5, 1.4) {structlog\\observability};
        \node[comp, minimum width=7.0cm] at (-3.8, -0.4) {6 agents (JD, Match, NLP, Vision, Scoring, Recruit)};
        \node[comp] at ( 1.8, -0.4) {Behavioural-signal\\fuser};
        \node[comp] at ( 5.5, -0.4) {PDF\\exporter};

        % --- Data tier ---
        \node[sideband=1.6cm] (datalbl) at (-9.0, -3.0)
            {Data\\tier};
        \node[band=1.6cm, fill=green!7] (data) at (0.5, -3.0) {};
        \node[comp, minimum width=2.4cm] at (-5.6, -3.0) {SQLite\\(\textit{hireflow.db})};
        \node[comp, minimum width=2.4cm] at (-2.8, -3.0) {Local file\\directory};
        \node[comp, minimum width=2.4cm] at ( 0.0, -3.0) {Knowledge graph\\(Graph RAG)};
        \node[comp, minimum width=2.4cm] at ( 2.8, -3.0) {Redis streams\\(A2A pub/sub)};
        \node[comp, minimum width=2.4cm] at ( 5.6, -3.0) {Hugging Face\\model cache};

        % Tier-to-tier links
        \draw[link] (-3.0, 3.2) -- (-3.0, 2.4);
        \draw[link] ( 3.0, 3.2) -- node[right, font=\scriptsize] {HTTP / WS} ( 3.0, 2.4);
        \draw[link] (-2.8, -1.4) -- (-2.8, -2.2);
        \draw[link] ( 2.8, -1.4) -- node[right, font=\scriptsize] {ORM / file I/O} ( 2.8, -2.2);
    \end{tikzpicture}%
    }}
    \caption{Three-tier logical architecture of HireFlow.}
    \label{fig:logical-arch}
\end{figure}

The presentation tier is a single-page application built with React
and styled with Tailwind. It hosts a Three.js avatar through React
Three Fiber and streams the candidate's microphone through WebRTC
to the backend. The business tier exposes a FastAPI gateway that
dispatches each call to the relevant agent through a LangGraph
state machine, with inter-agent traffic carried over Redis streams
that follow the \acs{A2A} contract. The data tier holds a SQLite
file for transactional state, a local file directory for binary
assets, an in-process talent knowledge graph queried by the JD
agent through Graph \acs{RAG}, and a shared Hugging Face cache for
model weights.

\subsection{Logical Architecture of the Multi-Agent System}

Inside the business tier, HireFlow follows a hub-and-spoke
multi-agent pattern. The \textbf{Recruit Agent} acts as the
orchestrator: it receives the recruiter's intent, decomposes it into
specialised subtasks and routes each subtask to one of the five
specialist agents. Communication between agents follows the
\textbf{\acs{A2A}} protocol over JSON-RPC and is materialised on top
of Redis streams, while each agent embeds its own models and tools
(the talent knowledge graph queried through Graph \acs{RAG}, vLLM,
the \acs{GNN} matcher, MediaPipe, Kokoro \acs{TTS}, Moonshine)
directly.

\begin{figure}[H]
    \centering
    \makebox[\textwidth][c]{\resizebox{1.10\textwidth}{!}{%
    \begin{tikzpicture}[
        font=\footnotesize,
        orch/.style={rectangle, draw=black!75, very thick, fill=orange!25,
            rounded corners=6pt, minimum width=4.6cm, minimum height=1.4cm,
            align=center, font=\footnotesize\bfseries, inner sep=4pt},
        agent/.style={rectangle, draw=black!70, thick, fill=blue!12,
            rounded corners=4pt, minimum width=2.8cm, minimum height=1.1cm,
            align=center, font=\footnotesize\bfseries, inner sep=3pt},
        tool/.style={rectangle, draw=black!60, thick, fill=green!10,
            rounded corners=2pt, minimum width=2.8cm, minimum height=1.3cm,
            align=center, font=\scriptsize, inner sep=2pt},
        layerband/.style={rectangle, draw=black!55, very thick, dashed,
            rounded corners=4pt, minimum width=16.6cm, inner sep=4pt},
        a2a/.style={thick, -{Latex[length=2.6mm]}, black!75},
        mcp/.style={thick, -{Latex[length=2.6mm]}, black!60, dashed}
    ]
        % Background bands for visual grouping (drawn first, behind agents).
        \node[layerband, minimum height=1.7cm] (agentband) at (0, 0.7)  {};
        \node[layerband, minimum height=2.0cm] (toolband)  at (0, -2.7) {};

        % Layer captions anchored on the top-left edge of each band so they
        % stay inside the visible figure area.
        \node[anchor=west, font=\scriptsize\bfseries\itshape, black!70,
              fill=white, inner sep=2pt]
              at ([xshift=8pt, yshift=-2pt]agentband.north west)
              {Agent layer};
        \node[anchor=west, font=\scriptsize\bfseries\itshape, black!70,
              fill=white, inner sep=2pt]
              at ([xshift=8pt, yshift=-2pt]toolband.north west)
              {Tool layer};

        % Orchestrator (centred, prominent)
        \node[orch] (recruit) at (0, 4.0) {Recruit Agent\\\textit{(orchestrator)}};
        \node[anchor=south, font=\scriptsize\bfseries\itshape, black!70,
              fill=white, inner sep=2pt]
              at ([yshift=4pt]recruit.north) {Orchestration};

        % 5 specialist agents — evenly spaced inside the agent band
        \node[agent] (jd)      at (-7.0, 0.7) {JD Agent};
        \node[agent] (match)   at (-3.5, 0.7) {Matching Agent};
        \node[agent] (nlp)     at ( 0.0, 0.7) {NLP Agent};
        \node[agent] (vision)  at ( 3.5, 0.7) {Vision Agent};
        \node[agent] (scoring) at ( 7.0, 0.7) {Scoring Agent};

        % Models / local tools — below their owning agent
        \node[tool] (jdllm) at (-7.0, -2.7) {Knowledge graph\\(Graph RAG)\\+ vLLM};
        \node[tool] (gnn)   at (-3.5, -2.7) {GNN matcher\\(TAPJFNN)};
        \node[tool] (vllm)  at ( 0.0, -2.7) {vLLM + Kokoro TTS\\+ Moonshine};
        \node[tool] (mp)    at ( 3.5, -2.7) {MediaPipe\\+ HSEmotion};
        \node[tool] (pdf)   at ( 7.0, -2.7) {Groq Cloud\\+ PDF export};

        % A2A links from orchestrator to each agent (fan-out)
        \draw[a2a] (recruit.south west) to[out=-150, in=90] (jd.north);
        \draw[a2a] (recruit.south west) to[out=-115, in=90] (match.north);
        \draw[a2a] (recruit.south)      --                  (nlp.north);
        \draw[a2a] (recruit.south east) to[out=-65,  in=90] (vision.north);
        \draw[a2a] (recruit.south east) to[out=-30,  in=90] (scoring.north);

        % Cross-agent A2A: NLP <-> Vision during an interview turn — routed
        % through the empty space between the agent and tool bands so the
        % label is clearly readable.
        \draw[a2a, <->] (nlp.south east) to[out=-30, in=-150]
            (vision.south west);
        \node[font=\scriptsize\itshape, fill=white, inner sep=2pt,
              draw=black!25, rounded corners=1.5pt]
              at (1.75, -1.1) {A2A (behavioural)};

        % Tool dependency links (agents to local models)
        \draw[mcp] (jd.south)      -- (jdllm.north);
        \draw[mcp] (match.south)   -- (gnn.north);
        \draw[mcp] (nlp.south)     -- (vllm.north);
        \draw[mcp] (vision.south)  -- (mp.north);
        \draw[mcp] (scoring.south) -- (pdf.north);

        % Compact legend (above the orchestrator, clear of all agents).
        \draw[black!50, thick, rounded corners=2pt, fill=white]
            (-1.6, 5.6) rectangle (1.6, 6.6);
        \node[anchor=west, font=\scriptsize\bfseries] at (-1.5, 6.45) {Legend};
        \draw[a2a] (-1.4, 6.00) -- (-0.7, 6.00);
        \node[anchor=west, font=\scriptsize] at (-0.65, 6.00) {A2A};
        \draw[mcp] ( 0.2, 6.00) -- ( 0.9, 6.00);
        \node[anchor=west, font=\scriptsize] at ( 0.95, 6.00) {uses};
    \end{tikzpicture}%
    }}
    \caption{Logical architecture of the HireFlow multi-agent system.}
    \label{fig:multiagent-arch}
\end{figure}

The shape of the diagram captures three design intentions. First,
\textbf{specialisation}: every agent owns a single concern --- JD
authoring, semantic matching, conversational \acs{NLP}, behavioural
vision, or scoring --- which keeps prompts compact and weights
swappable. Second, \textbf{co-located tooling}: each agent embeds
its own models and runs in-process, so the deployment stays a
single \acs{GPU} node with no external broker beyond Redis. Third,
\textbf{lateral communication}: the NLP and Vision agents talk
directly to each other during an interview turn, because the
orchestrator does not need to mediate the per-frame stream that
fuses speech with behavioural cues.

\clearpage
\section{Work Environment}

This section lists the hardware and software stack that supported
the implementation, the training experiments and the writing of the
present report.

\subsection{Hardware Environment}

Two machines were used throughout the internship: a developer
laptop for everyday coding and report writing, and a Lightning AI
Studio \acs{GPU} node for training experiments and inference
benchmarking. Table~\ref{tab:hardware} contrasts their
specifications.

\begin{table}[H]
    \centering
    {\renewcommand{\arraystretch}{1.1}%
    \resizebox{\textwidth}{!}{%
    \begin{tabular}{|l|p{5.5cm}|p{5.5cm}|}
        \hline
        \textbf{Component} & \textbf{Developer laptop} & \textbf{Lightning AI Studio} \\
        \hline
        CPU              & AMD Ryzen 5 7535HS, 6 cores  & 16 vCPU \\
        \hline
        GPU              & NVIDIA RTX 2050, 4~GB VRAM   & NVIDIA RTX 6000 Pro, 96~GB VRAM \\
        \hline
        RAM              & 16~GB                        & 32~GB \\
        \hline
        Storage          & 512~GB SSD                   & 200~GB persistent volume \\
        \hline
        Connectivity     & 100~Mbps                     & 1~Gbps \\
        \hline
        Role             & Coding, debugging, report    & Training, inference, demos \\
        \hline
    \end{tabular}%
    }}
    \caption{Hardware environment used during the internship.}
    \label{tab:hardware}
\end{table}

\clearpage
\subsection{Software Environment}

The software stack is open-source end to end. The choices favour
mature, well-documented projects with active communities so that the
platform remains maintainable beyond the duration of the internship.
Table~\ref{tab:software} groups the tools by responsibility.

\begin{table}[H]
    \centering
    {\renewcommand{\arraystretch}{1.05}%
    \resizebox{\textwidth}{!}{%
    \begin{tabular}{|l|p{11.0cm}|}
        \hline
        \textbf{Layer} & \textbf{Tools and libraries} \\
        \hline
        Operating system   & Ubuntu 22.04 LTS (server), Zorin OS (laptop) \\
        \hline
        Languages          & Python 3.11, TypeScript 5, SQL \\
        \hline
        Backend            & FastAPI, Uvicorn, Pydantic v2, sqlite3, structlog \\
        \hline
        AI/ML              & PyTorch 2, Hugging Face Transformers, PEFT/LoRA,
                             vLLM, LangGraph \\
        \hline
        Voice              & Kokoro TTS, Moonshine, Silero VAD \\
        \hline
        Vision             & MediaPipe, OpenCV, HSEmotion, Wav2Vec2 \\
        \hline
        Frontend           & React 19, Vite 7, TailwindCSS 4, Three.js, R3F, Drei \\
        \hline
        Persistence        & SQLite 3, Redis streams, local file directory \\
        \hline
        External APIs      & Groq Cloud (scoring LLM), Hugging Face Hub (weights) \\
        \hline
        IDE \& tooling     & Visual Studio Code, Git, GitHub \\
        \hline
        Documentation      & \LaTeX{} (TeXLive 2024), draw.io, TikZ \\
        \hline
    \end{tabular}%
    }}
    \caption{Software environment used during the internship.}
    \label{tab:software}
\end{table}

\section*{Conclusion}

This chapter translated HireFlow's high-level objectives into a
concrete engineering plan. The functional and non-functional
requirements were enumerated and exposed through a global use-case
diagram; a Scrum cadence with a fifteen-story product backlog and
a six-sprint roadmap was framed; and the physical, three-tier
logical and multi-agent architectures were detailed alongside the
hardware and software environment that supports them. The next
chapter reports on the realisation of these models and on the
experimental validation of each agent.
